\subsection{DERIVATIVE WITH RESPECT TO A PARAMETER OF AN INTEGRAL OVER A NON-COMPACT INTERVAL}

The set V having the same meaning as in prop. 5 of II, p.~74, suppose now that I is any interval in R, and that f is a continuous map from I × V into E; if the integral \(\mathbf{g}(\alpha) = \int_{\mathrm{I}} \mathbf{f}(t, \alpha) dt\) exists for all \(\alpha \in V\) and is a continuous function of \(\alpha\) , the function g need not have a derivative equal to \(\int_{\mathrm{I}} \mathbf{f}'_{\alpha}(t, \alpha_0) dt\) at the point \(\alpha_0\) , even if \(\mathbf{f}'_{\alpha}(x, \alpha)\)

converges uniformly to \(\mathbf{f}_{\alpha}^{\prime}(x,\alpha_{0})\) on every compact interval contained in I, and if the integral \(\int_{\mathrm{I}}\mathbf{f}_{\alpha}^{\prime}(t,\alpha)dt\) exists for all \(\alpha \in \mathrm{V}\) (cf.~II, p.~87, exerc. 3).

A sufficient condition for formula (12) (II, p.~74) to remain valid is given by the following proposition:

PROPOSITION 7. Let I be an arbitrary interval in \(\mathbf{R}\) , and \(\mathbf{f}\) a continuous function on \(\mathrm{I} \times \mathrm{V}\) . Suppose that:

\(1^{\circ}\) the partial derivative \(\mathbf{f}_{\alpha}^{\prime}(x,\alpha)\) exists for all \(x\in \mathrm{I}\) and all \(\alpha \in \mathrm{V}\) , and, for all \(\alpha \in \mathrm{V}\) , the map \(x\mapsto \mathbf{f}_{\alpha}^{\prime}(x,\alpha)\) is regulated on \(\mathrm{I}\) ;

\(2^{\circ}\) for all \(\alpha \in \mathrm{V}\) , \(\mathbf{f}_{\alpha}^{\prime}(x,\beta)\) converges uniformly to \(\mathbf{f}_{\alpha}^{\prime}(x,\alpha)\) on every compact interval contained in I, as \(\beta\) tends to \(\alpha\) ;

\(3^{\circ}\) the integral \(\int_{\mathrm{I}}\mathbf{f}_{\alpha}^{\prime}(t,\alpha)dt\) is uniformly convergent on \(\mathrm{V}\) ;

\(4^{\circ}\) the integral \(\int_{\mathrm{I}}\mathbf{f}(t,\alpha_0)dt\) is convergent.

In these circumstances the integral \(\mathbf{g}(\alpha)=\int_{\mathrm{I}}\mathbf{f}(t,\alpha)dt\) is uniformly convergent on V, and the function g admits at every point of V a derivative given by the formula

\[
\mathbf {g} ^ {\prime} (\alpha) = \int_ {\mathrm{I}} \mathbf {f} _ {\alpha} ^ {\prime} (t, \alpha) d t.\tag{14}
\]

The uniform convergence of \(\int_{\mathrm{I}} \mathbf{f}_{\alpha}'(t, \alpha) dt\) on V means that the function \(\alpha \mapsto \int_{\mathrm{J}} \mathbf{f}_{\alpha}'(t, \alpha) dt\) converges uniformly on V with respect to the filter of sections \(\Phi\) of the directed set \(\Re(\mathrm{I})\) of compact intervals J contained in I. Let us put \(\mathbf{u}_{\mathrm{J}}(\alpha) = \int_{\mathrm{J}} \mathbf{f}(t, \alpha) dt\) ; the hypotheses show that on the one hand \(\mathbf{u}_{\mathrm{J}}(\alpha_0)\) has a limit with respect to \(\Phi\) , and on the other hand, by virtue of prop. 5 of II, p.~74, that \(\mathbf{u}_{\mathrm{J}}'(\alpha) = \int_{\mathrm{J}} \mathbf{f}_{\alpha}'(t, \alpha) dt\) for all \(\alpha \in \mathrm{V}\) . We can therefore apply th. 1 of II, p.~52 to the functions \(\mathbf{u}_{\mathrm{J}}\) , the rôle of the set of indices being taken here by \(\Re(\mathrm{I})\) , and that of the filter on this set by the filter \(\Phi\) ; the proposition follows immediately.

Remarks. 1) Conditions \(1^{\circ}\) and \(2^{\circ}\) of prop. 7 are satisfied a fortiori when \(\mathbf{f}_{\alpha}^{\prime}(x,\alpha)\) is a continuous function of \((x,\alpha)\) on \(\mathrm{I} \times \mathrm{V}\) .

\begin{enumerate}
\def\labelenumi{\arabic{enumi})}
\setcounter{enumi}{1}
\tightlist
\item
  When, in an integral \(\int_{a(\alpha)}^{b(\alpha)}\mathbf{f}(t,\alpha)dt\) , the endpoints of the interval are finite functions of the parameter, the study of this integral as a function of \(\alpha\) can be related to that of an integral over {[}0, 1{]}; indeed, by the change of variable \(t = a(\alpha)(1 - u) + b(\alpha)u\) , one has
\end{enumerate}

\[
\int_ {a (\alpha)} ^ {b (\alpha)} \mathbf {f} (t, \alpha) d t = \int_ {0} ^ {1} \mathbf {f} (a (\alpha) (1 - u) + b (\alpha) u, \alpha) (b (\alpha) - a (\alpha)) d u.
\]

